% !TEX program = xelatex
% ============================================================
%  第 2 章 Convex Sets — 基于 template/main.tex 的视觉结构
%  正文与定理写法保持不变；页面为 A4 + ctex + xelatex
% ============================================================

\documentclass[11pt,a4paper]{article}

% ------------------------------------------------------------
%  字体与页面
% ------------------------------------------------------------
% 有 geometry 时 2.5cm 边距；精简 TeX 无该宏包时退化为标准 article 的 a4 版心
\IfFileExists{geometry.sty}{%
  \usepackage[margin=2.5cm]{geometry}
}{%
  \PackageWarning{geometry}{未找到 geometry 宏包，已使用默认 A4 页边距。建议安装
    宏包 `geometry' 以获得 2.5cm 边距。}%
}
\usepackage{ctex}
\usepackage{microtype}
\linespread{1.08}

% ------------------------------------------------------------
%  工具宏包
% ------------------------------------------------------------
\usepackage{amsmath,amssymb,amsthm}
\usepackage{mathtools}
\usepackage{graphicx}
\usepackage{booktabs}
\usepackage{enumitem}
\usepackage{tikz}
\usetikzlibrary{arrows,calc,positioning}
\usepackage[dvipsnames]{xcolor}
\usepackage[most]{tcolorbox}
\usepackage{titlesec}
\usepackage{setspace}
\setlength{\parskip}{0.35em}
\setlength{\parindent}{2em}
\setlist[itemize]{itemsep=0.2em, topsep=0.25em}
\setlist[enumerate]{itemsep=0.2em, topsep=0.25em}

% ------------------------------------------------------------
%  颜色与简易高亮（与模板一致，避免依赖 soul+中文）
% ------------------------------------------------------------
\definecolor{mydarkblue}{rgb}{0.0,0.08,0.45}
\definecolor{cardbg}{RGB}{245,245,247}
\definecolor{highlight}{HTML}{FCD66F}
\definecolor{warning}{HTML}{E57373}
\definecolor{info}{HTML}{4DB6AC}
\newcommand{\tlight}[1]{\textcolor{warning}{#1}}
\newcommand{\tinfo} [1]{\textcolor{info}{#1}}
\newcommand{\cbox}  [2]{\colorbox{#1}{\textcolor{black}{#2}}}
\newcommand{\bglight}[1]{\cbox{highlight}{#1}}
\newcommand{\bgwarn} [1]{\cbox{warning}{#1}}
\newcommand{\bginfo} [1]{\cbox{info}{#1}}

% ------------------------------------------------------------
%  超链接
% ------------------------------------------------------------
\usepackage{hyperref}
\hypersetup{
  colorlinks = true,
  linkcolor  = mydarkblue,
  citecolor  = mydarkblue,
  urlcolor   = mydarkblue
}

% ------------------------------------------------------------
%  章节标题（与模板一致：section 带 titlerule）
% ------------------------------------------------------------
\titleformat{\section}
  {\large\bfseries\color{mydarkblue}}{\thesection}{0.5em}{}[\titlerule]
\titlespacing*{\section}{0pt}{1.5ex}{1ex}

\titleformat{\subsection}
  {\normalsize\bfseries}{\thesubsection}{0.5em}{}
\titlespacing*{\subsection}{0pt}{1ex}{0.5ex}

% ------------------------------------------------------------
%  模板风格：紧凑顶栏 + 横向导航
% ------------------------------------------------------------
\makeatletter
\newcommand{\makeChapterTwoTitle}{%
  \begin{tcolorbox}[
    colback=mydarkblue!5, colframe=mydarkblue!10,
    arc=3pt, boxrule=0.5pt,
    left=5pt, right=5pt, top=5pt, bottom=5pt
  ]
    {\large\bfseries\@title}\hfill{\small\@author}\\[-2pt]
    {\tiny\color{gray}\@date\hfill Chapter 2 · Convex sets notes}
  \end{tcolorbox}
  \vspace{-0.6em}
}
\makeatother

\newtcolorbox{tabNavBoxSmall}[1][]{
  enhanced,
  colframe=white, colback=gray!5, coltext=black!80,
  fontupper=\sffamily\small\bfseries,
  arc=1.5mm, boxrule=0pt, bottomrule=2pt,
  left=2mm, right=2mm, top=1.5mm, bottom=1.5mm,
  #1
}

\newcommand{\horizontalTOCChTwo}{%
  \begin{center}
  \begin{tcbraster}[
    raster columns=3, raster width=\linewidth,
    raster column skip=1mm, raster equal height
  ]
    \begin{tabNavBoxSmall}
      \hyperref[sec:ch2-basic]{\ 优化与基本形式}
    \end{tabNavBoxSmall}
    \begin{tabNavBoxSmall}
      \hyperref[sec:ch2-affine]{\ Affine sets 仿射集}
    \end{tabNavBoxSmall}
    \begin{tabNavBoxSmall}
      \hyperref[sec:ch2-convex]{\ Convex sets 凸集}
    \end{tabNavBoxSmall}
    \begin{tabNavBoxSmall}
      \hyperref[sec:ch2-cone]{\ Convex cones 凸锥}
    \end{tabNavBoxSmall}
    \begin{tabNavBoxSmall}
      \hyperref[sec:ch2-linear-fractional]{\ 线性分数函数}
    \end{tabNavBoxSmall}
    \begin{tabNavBoxSmall}
      \hyperref[sec:ch2-separation-support]{\ 分离与支撑超平面}
    \end{tabNavBoxSmall}
  \end{tcbraster}
  \vspace{-2.5pt}%
  \noindent\textcolor{gray!30}{\rule{\linewidth}{1pt}}
  \end{center}
  \vspace{1.2em}
}

% ------------------------------------------------------------
%  定理环境：与 template/main.tex 同款 newtcbtheorem
%  用法：\begin{definition}{标题}{编号} ... \end{definition}
% ------------------------------------------------------------
\tcbset{
  mytheo/.style={
    enhanced, breakable,
    colback=#1!5!white, colframe=#1!70!black,
    coltitle=white, fonttitle=\bfseries\small,
    boxrule=0.6pt, arc=3pt,
    left=5pt, right=5pt, top=10pt, bottom=5pt,
    attach boxed title to top left={yshift*=-\tcboxedtitleheight/2, xshift=2mm},
    boxed title style={
      colback=#1!80!black, colframe=#1!80!black,
      boxrule=0pt, arc=2pt
    }
  }
}
\newtcbtheorem[auto counter]            {theorem}    {定理}{mytheo=WildStrawberry}{thm}
\newtcbtheorem[use counter from=theorem]{definition} {定义}{mytheo=NavyBlue}{def}
\newtcbtheorem[use counter from=theorem]{example}    {例子}{mytheo=YellowOrange}{ex}
\newtcbtheorem[use counter from=theorem]{remark}     {注解}{mytheo=Periwinkle}{rem}

% 证明：轻量底纹
\tcolorboxenvironment{proof}{
  blanker, breakable, left=4pt, right=4pt, top=3pt, bottom=3pt,
  borderline west={2pt}{0pt}{gray!50}
}

\title{第 2 章\quad Convex Sets 笔记}
\author{}
\date{\today}

\begin{document}

\makeChapterTwoTitle
\horizontalTOCChTwo

\begin{center}
  \includegraphics[width=\linewidth]{\detokenize{ChatGPT Image 2026年4月26日 19_03_54.png}}
\end{center}
\vspace{0.5em}

\section{优化问题基本形式}\label{sec:ch2-basic}
\[
\begin{aligned}
\text{minimize} \quad & f_0(x) \\
\text{subject to} \quad & f_i(x) \leq b_i,\quad i=1,2,\dots,m \\
& x = \begin{bmatrix} x_1 \\ x_2 \\ \vdots \\ x_n \end{bmatrix} \in \mathbb{R}^n
\end{aligned}
\]

Convex / non-convex $\quad \to \quad$ \textbf{简单 / 不简单}

\textbf{重点章节}：2、3、4、5、9、10

\section{Affine Sets 仿射集}\label{sec:ch2-affine}

\subsection{直线（Line）}
设 $x_1, x_2 \in \mathbb{R}^n$，$\theta \in \mathbb{R}$，则
\[
y = \theta x_1 + (1-\theta)x_2 = x_2 + \theta(x_1 - x_2)
\]

\subsection{线段（Line Segment）}
\[
y = \theta x_1 + (1-\theta)x_2,\quad 0 \leq \theta \leq 1
\]

\begin{definition}{仿射集}{affine-set}
集合 $C \subseteq \mathbb{R}^n$ 称为\textbf{仿射集}（Affine Set），当且仅当对任意 $x_1, x_2 \in C$ 和任意 $\theta \in \mathbb{R}$，都有
\[
\theta x_1 + (1-\theta)x_2 \in C\,.
\]
\end{definition}

\begin{definition}{仿射组合}{affine-combination}
设 $x_1, x_2, \dots, x_k \in C$，$\theta_1, \theta_2, \dots, \theta_k \in \mathbb{R}$ 且满足
\[
\sum_{i=1}^k \theta_i = 1
\]
则称
\[
\sum_{i=1}^k \theta_i x_i
\]
为 $x_1,x_2,\dots,x_k$ 的一个\textbf{仿射组合}（Affine Combination）。
\end{definition}

\begin{theorem}{仿射组合封闭性}{affine-closure}
若 $C$ 是仿射集，则 $C$ 中任意有限个点的仿射组合仍属于 $C$。
\end{theorem}

\subsection{仿射组合的封闭性（三点示例）}
设 $x_1, x_2, x_3 \in C$，$\theta_1, \theta_2, \theta_3 \in \mathbb{R}$ 且 $\theta_1 + \theta_2 + \theta_3 = 1$，则
\[
\theta_1 x_1 + \theta_2 x_2 + \theta_3 x_3 \in C\,.
\]

\subsection{仿射集的平移表示}
设 $C$ 为仿射集，$x_0 \in C$，定义
\[
V = C - x_0 = \{ x - x_0 \mid x \in C \}\,.
\]
则 $V$ 是线性子空间（即对任意 $u_1, u_2 \in V$，$\alpha, \beta \in \mathbb{R}$，都有 $\alpha u_1 + \beta u_2 \in V$）。

\begin{proof}[简证]
任取 $u_1, u_2 \in V$，则存在 $x_1, x_2 \in C$ 使得 $u_1 = x_1 - x_0$，$u_2 = x_2 - x_0$。于是
\[
\alpha u_1 + \beta u_2 = [\alpha x_1 + \beta x_2 + (1-\alpha-\beta)x_0] - x_0\,.
\]
因为 $C$ 是仿射集，上式括号内向量属于 $C$，故 $\alpha u_1 + \beta u_2 \in V$。
\end{proof}

\subsection{线性方程组的解集是仿射集}
\begin{example}{线性方程组解集是仿射集}{affine-linear-system}
设
\[
C = \{ x \in \mathbb{R}^n \mid Ax = b \},\quad A \in \mathbb{R}^{m \times n},\ b \in \mathbb{R}^m\,.
\]
则 $C$ 是仿射集。
\end{example}

\begin{proof}
任取 $x_1, x_2 \in C$，$\theta \in \mathbb{R}$，则
\[
A(\theta x_1 + (1-\theta)x_2) = \theta b + (1-\theta)b = b\,.
\]
因此 $\theta x_1 + (1-\theta)x_2 \in C$。
\end{proof}

\begin{remark}{解集的平移表示}{affine-translation}
若 $x_0 \in C$（即 $A x_0 = b$），则
\[
V = C - x_0 = \{ y \mid A y = 0 \}\,,
\]
即 $V$ 是齐次线性方程组的解空间（零空间）。
\end{remark}

\subsection{仿射包（Affine Hull）——最小仿射集}
\begin{definition}{仿射包}{affine-hull}
任意集合 $C \subseteq \mathbb{R}^n$ 的\textbf{仿射包}（Affine Hull）定义为
\[
\operatorname{aff} C = \left\{ \sum_{i=1}^k \theta_i x_i \;\middle|\; x_i \in C,\ \theta_i \in \mathbb{R},\ \sum_{i=1}^k \theta_i = 1 \right\}\,.
\]
\end{definition}

\begin{theorem}{仿射包最小性}{affine-hull-minimal}
$\operatorname{aff} C$ 是包含 $C$ 的最小仿射集：即若 $S$ 是任意仿射集且 $C \subseteq S$，则 $\operatorname{aff} C \subseteq S$。
\end{theorem}

\section{Convex Sets 凸集}\label{sec:ch2-convex}

\begin{definition}{凸集}{convex-set}
集合 $C \subseteq \mathbb{R}^n$ 称为\textbf{凸集}（Convex Set），当且仅当对任意 $x_1, x_2 \in C$ 和任意 $\theta \in [0,1]$，都有
\[
\theta x_1 + (1-\theta)x_2 \in C\,.
\]
\end{definition}

\begin{remark}{仿射集与凸集关系}{affine-vs-convex}
仿射集是凸集的特例（因为仿射集对任意 $\theta \in \mathbb{R}$ 封闭，而凸集只需 $\theta \in [0,1]$）。
\end{remark}

\begin{definition}{凸组合}{convex-combination}
设 $x_1, x_2, \dots, x_k \in C$，$\theta_1, \theta_2, \dots, \theta_k \in \mathbb{R}$ 满足
\[
\theta_i \geq 0,\quad \sum_{i=1}^k \theta_i = 1\,,
\]
则称 $\sum_{i=1}^k \theta_i x_i$ 为 $x_1,\dots,x_k$ 的一个\textbf{凸组合}（Convex Combination）。
\end{definition}

\begin{theorem}{凸组合封闭性判别}{convex-combination-criterion}
$C$ 是凸集当且仅当 $C$ 中任意有限个点的凸组合仍属于 $C$。
\end{theorem}

\subsection{凸包（Convex Hull）——最小凸集}
\begin{definition}{凸包}{convex-hull}
任意集合 $C \subseteq \mathbb{R}^n$ 的\textbf{凸包}（Convex Hull）定义为
\[
\operatorname{conv} C = \left\{ \sum_{i=1}^k \theta_i x_i \;\middle|\; x_i \in C,\ \theta_i \geq 0,\ \sum_{i=1}^k \theta_i = 1 \right\}\,.
\]
\end{definition}

\begin{theorem}{凸包最小性}{convex-hull-minimal}
$\operatorname{conv} C$ 是包含 $C$ 的最小凸集。
\end{theorem}

\section{Convex Cones 凸锥}\label{sec:ch2-cone}

\begin{definition}{凸锥}{convex-cone}
集合 $C \subseteq \mathbb{R}^n$ 称为\textbf{凸锥}（Convex Cone），当且仅当
\[
x \in C,\ \theta \geq 0 \quad \Longrightarrow \quad \theta x \in C\,.
\]
（等价地：对任意 $x_1, x_2 \in C$ 和 $\theta_1, \theta_2 \geq 0$，都有 $\theta_1 x_1 + \theta_2 x_2 \in C$）
\end{definition}

\begin{definition}{凸锥组合}{convex-cone-combination}
设 $x_1,\dots,x_k \in C$，$\theta_1,\dots,\theta_k \geq 0$，则
\[
\sum_{i=1}^k \theta_i x_i
\]
称为凸锥组合。
\end{definition}

\begin{remark}{凸锥包记号}{convex-cone-notation}
凸锥包（最小凸锥）是包含 $C$ 的最小凸锥，记为 $\operatorname{cone} C$ 或 $\operatorname{convcone} C$。
\end{remark}

\section{几种重要的凸集}

\begin{definition}{重要凸集示例}{important-convex-sets}
以下均为凸集（部分同时也是仿射集）：
\begin{itemize}[leftmargin=*, itemsep=0.5em]
\item $\mathbb{R}^n$ 整个空间
\item $\mathbb{R}^n$ 的任意子空间（subspace）
\item 任意直线（line）
\item 任意线段（line segment）（凸集，但一般不是仿射集）
\item 任意射线（ray）：$\{x_0 + \theta v \mid \theta \geq 0\}$，其中 $x_0 \in \mathbb{R}^n$，$v \in \mathbb{R}^n$
\end{itemize}
任意子空间既是凸集也是仿射集。
\end{definition}

\section{超平面与半空间}

\subsection{超平面（Hyperplane）}
\[
\{ x \in \mathbb{R}^n \mid a^T x = b \},\quad a \in \mathbb{R}^n \setminus \{0\},\ b \in \mathbb{R}
\]

\subsection{半空间（Halfspace）}
\[
\{ x \in \mathbb{R}^n \mid a^T x \leq b \}
\]
（同理 $\{ x \mid a^T x \geq b \}$ 也是半空间）

\begin{remark}{超平面与半空间性质}{hyperplane-halfspace}
超平面是仿射集，半空间是凸集（但不是仿射集）。
\end{remark}

\section{球与椭球}

\subsection{欧氏球（Euclidean Ball）}
\[
\mathcal{B}(x_c, r) = \{ x \in \mathbb{R}^n \mid \|x - x_c\|_2 \leq r \}
= \left\{ x \;\middle|\; \sqrt{(x - x_c)^T (x - x_c)} \leq r \right\}
\]

\begin{proof}[凸性证明]
任取 $x_1, x_2 \in \mathcal{B}(x_c, r)$，$\theta \in [0,1]$，则
\begin{align*}
\|\theta x_1 + (1-\theta)x_2 - x_c\|_2
&= \|\theta(x_1 - x_c) + (1-\theta)(x_2 - x_c)\|_2 \\
&\leq \theta \|x_1 - x_c\|_2 + (1-\theta)\|x_2 - x_c\|_2 \\
&\leq \theta r + (1-\theta)r = r.
\end{align*}
故 $\theta x_1 + (1-\theta)x_2 \in \mathcal{B}(x_c, r)$。
\end{proof}

\subsection{椭球（Ellipsoid）}
\[
\mathcal{E} = \left\{ x \in \mathbb{R}^n \;\middle|\; (x - x_c)^T P^{-1} (x - x_c) \leq 1 \right\},\quad P \in \mathbf{S}_{++}^n
\]
（$\mathbf{S}_{++}^n$ 表示正定矩阵）

特别地，当 $P = r^2 I_n$ 时，椭球退化为欧氏球。

\begin{remark}{奇异值补充}{ellipsoid-svd}
若 $P = A^T A$（$A$ 可逆），则可通过奇异值分解表示椭球形状。
\end{remark}

\subsection{正定矩阵（Positive Definite）}
设 $A \in \mathbb{R}^{n \times n}$，若对任意 $x \neq 0$ 都有 $x^T A x > 0$，则称 $A$ 为\textbf{正定矩阵}。

\begin{enumerate}
\item 其特征值 $\lambda_i(A) > 0$（全正）
\item 存在 Cholesky 分解：$A = LL^T$（$L$ 为下三角矩阵）
\item 函数 $f(x) = x^T A x$ 是严格凸函数（strictly convex）
\end{enumerate}

\section{多面体（Polyhedron）}

\begin{definition}{多面体}{polyhedron}
集合 $P \subseteq \mathbb{R}^n$ 称为\textbf{多面体}（Polyhedron），若它可表示为有限个半空间与超平面的交集：
\[
P = \{ x \in \mathbb{R}^n \mid a_j^T x \leq b_j,\ j=1,\dots,m;\quad c_j^T x = d_j,\ j=1,\dots,p \}
\]
其中 $a_j, c_j \in \mathbb{R}^n$，$b_j, d_j \in \mathbb{R}$。
\end{definition}

等价地，$P = \{ x \mid Ax \leq b,\ Cx = d \}$，其中
\[
A = \begin{bmatrix} a_1^T \\ \vdots \\ a_m^T \end{bmatrix},\quad
C = \begin{bmatrix} c_1^T \\ \vdots \\ c_p^T \end{bmatrix}.
\]

\subsection{单纯形（Simplex）}
在 $\mathbb{R}^n$ 中选取 $k+1$ 个仿射无关的点 $v_0, v_1, \dots, v_k \in \mathbb{R}^n$，则它们的\textbf{单纯形}定义为
\[
C = \operatorname{conv}\{v_0, v_1, \dots, v_k\} = \left\{ \sum_{i=0}^k \theta_i v_i \;\middle|\; \theta_i \geq 0,\ \sum_{i=0}^k \theta_i = 1 \right\}.
\]

\begin{theorem}{单纯形与多面体}{simplex-is-polyhedron}
单纯形一定是多面体。
\end{theorem}

\begin{proof}
设 $x \in C$，则存在 $\theta_i \geq 0$，$\sum \theta_i = 1$ 使得 $x = \sum_{i=0}^k \theta_i v_i$。
令 $y = x - v_0 = \sum_{i=1}^k \theta_i (v_i - v_0)$，则
\[
y = B \theta',\quad \theta' = (\theta_1, \dots, \theta_k)^T,\quad B = [v_1 - v_0,\ \dots,\ v_k - v_0] \in \mathbb{R}^{n \times k}.
\]
因为 $v_0,\dots,v_k$ 仿射无关，所以 $\operatorname{rank}(B) = k$。存在矩阵 $A_1, A_2$ 使得
\[
A_1 B = I_k,\quad A_2 B = 0.
\]
从而
\[
A_1 y = \theta' \geq 0,\quad A_2 y = 0.
\]
代回 $x = y + v_0$ 可得
\[
A_1(x - v_0) \geq 0,\quad A_2(x - v_0) = 0,
\]
即 $x$ 满足一组线性不等式与等式约束，因此 $C$ 是多面体。
\end{proof}

\section{保持凸性的运算（Operations that preserve convexity）}

\subsection{交集（Intersection）}
若 $S_1, S_2$ 为凸集，则 $S_1 \cap S_2$ 为凸集。
更一般地，若对任意 $\alpha \in A$，$S_\alpha$ 均为凸集，则
\[
\bigcap_{\alpha \in A} S_\alpha
\]
仍为凸集。

\section{仿射函数（Affine Functions）}

\begin{definition}{仿射函数}{affine-function}
函数 $f: \mathbb{R}^n \to \mathbb{R}^m$ 称为\textbf{仿射函数}（Affine），当且仅当存在 $A \in \mathbb{R}^{m \times n}$ 和 $b \in \mathbb{R}^m$ 使得
\[
f(x) = Ax + b.
\]
\end{definition}

\begin{theorem}{仿射像保持凸性}{affine-forward-image}
设 $S \subseteq \mathbb{R}^n$ 为凸集，$f: \mathbb{R}^n \to \mathbb{R}^m$ 为仿射函数，则像集
\[
f(S) = \{ f(x) \mid x \in S \}
\]
为凸集。
\end{theorem}

\begin{theorem}{逆像（Inverse image）}{affine-inverse-image}
设 $g: \mathbb{R}^k \to \mathbb{R}^n$ 为仿射函数，$S \subseteq \mathbb{R}^n$ 为凸集，则逆像
\[
g^{-1}(S) = \{ x \in \mathbb{R}^k \mid g(x) \in S \}
\]
为凸集。
\end{theorem}

\section{缩放与移位（Scaling and Translation）}

缩放（scaling）和移位（translation）运算保持凸性：

\begin{itemize}
\item 若 $S$ 为凸集，$\lambda \geq 0$，则 $\lambda S = \{\lambda x \mid x \in S\}$ 为凸集；
\item 若 $S$ 为凸集，$y \in \mathbb{R}^n$，则 $S + y = \{x + y \mid x \in S\}$ 为凸集。
\end{itemize}

\begin{example}{凸集的和与笛卡尔积}{sum-cartesian-convex}
设 $S_1, S_2$ 为凸集，则
\[
S_1 + S_2 = \{ x + y \mid x \in S_1,\ y \in S_2 \}
\]
和
\[
S_1 \times S_2 = \{ (x,y) \mid x \in S_1,\ y \in S_2 \}
\]
均为凸集。
\end{example}

\section{半正定、LMI 与椭球（补充）}

\subsection{线性矩阵不等式（Linear Matrix Inequality, LMI）}
设 $A_1, A_2, \dots, A_n, B \in \mathbf{S}^m$（$m$ 阶对称矩阵），定义仿射映射
\[
\mathcal{A}(x) = \sum_{i=1}^n x_i A_i\,.
\]
则\textbf{线性矩阵不等式}（LMI）的解集为
\[
\{ x \in \mathbb{R}^n \mid \mathcal{A}(x) \preceq B \}
\]
（其中 $\preceq$ 表示 $B - \mathcal{A}(x)$ 是半正定矩阵）。

\begin{theorem}{LMI 的解集是凸集}{lmi-solution-convex}
LMI 的解集是凸集。
\end{theorem}

\begin{proof}
定义仿射变换 $f(x) = B - \mathcal{A}(x)$，则
\[
\{ x \mid \mathcal{A}(x) \preceq B \} = \{ x \mid f(x) \succeq 0 \} = f^{-1}(\mathbf{S}_+^m)\,,
\]
其中 $\mathbf{S}_+^m$ 是半正定锥，是凸集。因为仿射函数的逆像保持凸性，所以该解集为凸集。
\end{proof}

\subsection{椭球的另一种证明（通过单位球的仿射映射）}
单位球是凸集：
\[
\{ u \in \mathbb{R}^n \mid \|u\|_2 \leq 1 \}\,.
\]
考虑严格仿射映射
\[
f(u) = P^{1/2} u + x_c\,,
\]
其中 $P \in \mathbf{S}_{++}^n$（正定矩阵）。则椭球可表示为
\[
\mathcal{E} = f(\{ u \mid \|u\|_2 \leq 1 \})\,.
\]
因为仿射函数像集保持凸性，故椭球是凸集。

等价地，若 $x = P^{1/2} u + x_c$，则 $u = P^{-1/2}(x - x_c)$，代入单位球不等式得
\[
(x - x_c)^T P^{-1} (x - x_c) \leq 1\,.
\]

\section{透视函数（Perspective Function）}

\begin{definition}{透视函数}{perspective-function}
透视函数 $P: \mathbb{R}^{n+1} \to \mathbb{R}^n$ 定义为
\[
P(x,t) = \frac{x}{t}\,,
\]
定义域为 $\operatorname{dom} P = \mathbb{R}^n \times \mathbb{R}_{++}$（$t > 0$）。
\end{definition}

\begin{theorem}{透视函数保持（像的）凸性}{perspective-preserves-image-convexity}
透视函数保持凸性：若 $C \subseteq \operatorname{dom} P$ 是凸集，则 $P(C)$ 是凸集。
\end{theorem}

\begin{proof}
任取 $x,y \in C$，$x = (\tilde{x}, x_{n+1})$，$y = (\tilde{y}, y_{n+1})$，$x_{n+1}>0$，$y_{n+1}>0$，$\theta \in [0,1]$。
考虑线段上的点：
\[
z = \theta x + (1-\theta)y = \bigl( \theta \tilde{x} + (1-\theta)\tilde{y},\ \theta x_{n+1} + (1-\theta)y_{n+1} \bigr)\,.
\]
则
\[
P(z) = \frac{\theta \tilde{x} + (1-\theta)\tilde{y}}{\theta x_{n+1} + (1-\theta)y_{n+1}}
= \frac{ \theta x_{n+1} \dfrac{\tilde{x}}{x_{n+1}} + (1-\theta)y_{n+1} \dfrac{\tilde{y}}{y_{n+1}} }{ \theta x_{n+1} + (1-\theta)y_{n+1} }
= \mu P(x) + (1-\mu) P(y)\,,
\]
其中
\[
\mu = \frac{\theta x_{n+1}}{\theta x_{n+1} + (1-\theta)y_{n+1}} \in [0,1]\,.
\]
因此 $P(z)$ 是 $P(x)$ 与 $P(y)$ 的凸组合，故 $P(C)$ 为凸集。
\end{proof}

\section{线性分数函数（Linear Fractional Function）}\label{sec:ch2-linear-fractional}

\begin{definition}{线性分数函数}{linear-fractional-function}
设 $g: \mathbb{R}^n \to \mathbb{R}^{m+1}$ 为仿射函数：
\[
g(x) = \begin{bmatrix} A \\ c^T \end{bmatrix} x + \begin{bmatrix} b \\ d \end{bmatrix},\quad A \in \mathbb{R}^{m \times n},\ c \in \mathbb{R}^n,\ b \in \mathbb{R}^m,\ d \in \mathbb{R}\,.
\]
则线性分数函数 $f: \mathbb{R}^n \to \mathbb{R}^m$ 定义为透视函数与仿射函数的复合：
\[
f(x) = P(g(x)) = \frac{Ax + b}{c^T x + d}\,,
\]
定义域为
\[
\operatorname{dom} f = \{ x \in \mathbb{R}^n \mid c^T x + d > 0 \}\,.
\]
\end{definition}

\begin{theorem}{线性分数函数保持（像的）凸性}{linear-fractional-preserves-image-convexity}
线性分数函数是凸集到凸集的映射（即若 $C \subseteq \operatorname{dom} f$ 是凸集，则 $f(C)$ 是凸集）。
\end{theorem}

\begin{example}{概率论中的应用}{linear-fractional-probability-example}
设 $U,V$ 是两个离散随机变量，取值分别为 $\{1,2,\dots,n\}$ 和 $\{1,2,\dots,m\}$。
联合概率 $P_{ij} = P(U=i,V=j)$，条件概率 $f_{ij} = P(U=i \mid V=j)$。
则
\[
f_{ij} = \frac{P_{ij}}{\sum_{k=1}^n P_{kj}} = \frac{P_{ij}}{P(V=j)}\,.
\]
这正是线性分数函数的形式（分子为仿射映射，分母为线性形式）。
\end{example}

\section{分离超平面与支撑超平面}\label{sec:ch2-separation-support}

\subsection{分离超平面（Separating Hyperplane）}
\begin{theorem}{分离超平面定理}{separating-hyperplane-theorem}
设 $C,D \subseteq \mathbb{R}^n$ 是两个非空凸集，且 $C \cap D = \emptyset$，则存在 $a \in \mathbb{R}^n \setminus \{0\}$ 和 $b \in \mathbb{R}$ 使得
\[
a^T x \leq b \quad \forall x \in C\,,
\]
\[
a^T x \geq b \quad \forall x \in D\,.
\]
\end{theorem}

\subsection{支撑超平面（Supporting Hyperplane）}
\begin{theorem}{支撑超平面}{supporting-hyperplane-theorem}
设 $C \subseteq \mathbb{R}^n$ 是非空凸集，$x_0 \in \partial C$（$x_0$ 在 $C$ 的边界上），则存在非零向量 $a \in \mathbb{R}^n$ 使得
\[
a^T (x - x_0) \leq 0 \quad \forall x \in C\,.
\]
（即超平面 $\{ x \mid a^T (x - x_0) = 0 \}$ 在 $x_0$ 处支撑 $C$）
\end{theorem}

\begin{remark}{分离与支撑超平面的作用}{separating-supporting-hyperplane-remark}
分离超平面用于证明两个不相交凸集的存在性，支撑超平面则描述凸集在边界点的局部性质。两者在凸优化中是核心几何工具。
\end{remark}

\end{document}
